\subsection{A positional class which lifts once and fails at the next norm}

\begin{statement}{E.62}{Proved}{The complete local norm tower at seventeen and an actual second-step obstruction}
Use the genuine arithmetic quotient of E.38
and E.46,
\[
 \begin{gathered}
 \Lambda=\mathbb Z_3[[T-1]],\quad
 J=(3,T-1)\Lambda,\quad
 H=3\mathbb Z_3\ltimes J,\\
 (a,v)(b,w)=(a+b,v+T^a w),\qquad
 g(a,v)=\bigl(\chi(g)a,\chi(g)v(T^{\chi(g)})\bigr).
 \end{gathered}
 \tag{E62.1}
\]
For $n\ge1$ put
\[
 \begin{gathered}
 \Lambda_n=\Lambda/(T^{3^n}-1),\qquad
 J_n=\ker(\Lambda_n\xrightarrow{\rm aug}\mathbb F_3),\\
 H_n=H/\bigl(0\times(T^{3^n}-1)\Lambda\bigr)
                  =3\mathbb Z_3\ltimes J_n.
 \end{gathered}
 \tag{E62.2}
\]
Fix the actual first cocycle
$A=3\kappa_3(17)$.
Under the Kummer--Shapiro comparison of E.46,
the positional algebras and their transition
maps are exactly
\[
 \begin{gathered}
 R_n=\mathbb Q[X]/(X^{3^n}-17^3)
       =F_{n-1}\times L_{n-1},\\
 F_d=\mathbb Q(\alpha_d),\quad
       \alpha_d^{3^d}=17,\qquad
 L_d=K(\beta_d),\quad
       \beta_d^{3^d}=17\zeta_3,\quad
       K=\mathbb Q(\zeta_3),\\
 \alpha_{d+1}^3=\alpha_d,\quad
 \beta_{d+1}^3=\beta_d,\qquad
 \alpha_0=17,\quad\beta_0=17\zeta_3 .
 \end{gathered}
 \tag{E62.3}
\]
The transition on section coordinates is the
actual norm in these two field factors.

There is a global section class of $H_1$
over the fixed first cocycle whose positional
Kummer coordinates are
\[
                         (1,\zeta_3)\in
               \mathbb Q^{\times,\wedge_3}
                     \times K^{\times,\wedge_3}.
 \tag{E62.4}
\]
It lifts to a section class of $H_2$.
It has no lift to a section class of $H_3$,
and hence no lift to a section class of $H$.
This failure is already detected by a finite
residue-field quotient at seventeen.

The entire local section tower over
$\mathbb Q_{17}$ can be evaluated.
Choose the primitive ninth root
$\iota\in K_{17}=\mathbb Q_{17}(\zeta_3)$
whose residue is $2+8\bar\zeta_3$; thus
$\iota^3=\zeta_3$.
Put
\[
 \mathcal L=\{(e,f)\in\mathbb Z_3^2:
                         e+2f\in3\mathbb Z_3\}.
 \tag{E62.5}
\]
For every $n\ge1$, the full local fiber of
section conjugacy classes of $H_n$ over
$[A]$ is
\[
 \mathcal L\oplus\mathbb Z/9,\qquad
 (e,f,t)\longleftrightarrow
             (\alpha_{n-1}^{\,e},
                         \beta_{n-1}^{\,f}\iota^t).
 \tag{E62.6}
\]
Here the displayed field elements denote
their completions for multiplication by
$3^m$. The transition is precisely
\[
                  (e,f,t)\longmapsto(e,f,3t).
 \tag{E62.7}
\]
Consequently the full local fiber for $H$
is the lattice $\mathcal L$, and its
projection to the first positional level is
\[
           (e,f)\longmapsto
                   (17^e,(17\zeta_3)^f).
 \tag{E62.8}
\]
These are homeomorphisms for the natural
continuous-cohomology and three-adic
topologies. The common toric quotient of
E.46 is $(a,v)\mapsto(a/3,\operatorname{aug}(v)/3)$.
On these local classes it is the explicit map
\[
       (e,f)\longmapsto
       \left(\kappa_3(17),
                  \frac{e+2f}{3}\kappa_3(17)\right).
 \tag{E62.9}
\]
It is onto the entire second local toric
coordinate. Every such toric fiber is a copy
of $\mathbb Z_3$, parametrized by $f$, with
$e=3b-2f$ for second coordinate
$b\kappa_3(17)$.

The global class (E62.4) restricts to
$(e,f,t)=(0,0,3)$ in (E62.6).
Its first lifts have $t=1,4,7$;
none has a further lift.
Its total norm to $\mathbb Q$ is one,
so its second toric class is zero.
The positional class $(1,1)$ has the
same toric pair and does extend:
the literal global section is
$g\mapsto(A(g),0,g)$ of $H\rtimes G_{\mathbb Q}$.
Thus the failed positional lift and the
surjective toric map are related by the
actual maps (E62.6)--(E62.9).
The manuscript anchor is \S51.16, folios
733v--736v, scans~172--178, formulas~(371)--(379),
Part~E(c),(d).
\end{statement}

\begin{proof}
\emph{The exact section coordinates.}
For a section with first cocycle $A$, its
second coordinate is a cocycle for the
actual module action
\[
       g\cdot v=\chi(g)T^{A(g)}
                                 v(T^{\chi(g)}).
 \tag{E62.10}
\]
This action preserves $J_n$.
The Kummer--Shapiro map of E.46 identifies
$H^1(G_{\mathbb Q},\Lambda_{A,n})$
with $R_n^{\times,\wedge_3}$.
The augmentation is the total norm,
and the transition is induced by
root cubing, hence by the field norms
in (E62.3).
The two polynomial factors there are
$X^{3^{n-1}}-17$ and its complementary
factor. The first is Eisenstein at
seventeen; the complementary factor
over $K$ is
$X^{3^{n-1}}-17\zeta_3$, also
Eisenstein at the unique prime over
seventeen. These facts prove all
degrees and the stated field
embeddings.

Both globally and over $\mathbb Q_{17}$,
$H^0(\mathbb F_3(1))=0$.
The exact sequence
$0\to J_{A,n}\to\Lambda_{A,n}
\to\mathbb F_3(1)\to0$ therefore
identifies $H^1(J_{A,n})$ with the
kernel of the total norm modulo cubes.
Moreover $H^0(\mathbb Z_3(1))=0$,
so fixing a representative $A$ of its
first class leaves no first-coordinate
stabilizer. Thus these are the full
section conjugacy classes over $[A]$,
with the actual geometric conjugations
accounted for.

\emph{The literal first lift.}
Use the field $L_1$ and elements of E.57:
\[
 \beta_1^3=17\zeta_3,\qquad
 \pi=2-\zeta_3,\qquad
 g_1=-2-3\zeta_3-\zeta_3\beta_1,\qquad
 w=(g_1/\pi)^2 .
 \tag{E62.11}
\]
The actual norm is
$N_{L_1/K}(w)=\zeta_3$.
Hence $(1,w)\in R_2^{\times,\wedge_3}$
maps to $(1,\zeta_3)$.
The total norm of each is one,
so the preceding exact sequence places
both in the asserted $J$-cohomology
classes. This proves the global first
lift with an explicit element.
The identity pair at every positional
level is represented by the zero
second cocycle. Since $A$ is a genuine
$\mathbb Z_3(1)$-cocycle, the group law
in (E62.1) verifies directly that
$g\mapsto(A(g),0,g)$ is a full
section. Its common toric pair is
$(\kappa_3(17),0)$, which is also
the toric pair of (E62.4).

\emph{The local unit norm and the obstruction.}
The prime seventeen is inert in $K$.
Its completion has residue field
$\mathbb F_{289}$, whose unit group
has order $288=32\cdot9$.
The local extensions
\[
 F_{d,17}=\mathbb Q_{17}(\alpha_d),
 \qquad L_{d,17}=K_{17}(\beta_d)
 \tag{E62.12}
\]
are totally ramified of degree $3^d$,
by their Eisenstein polynomials.
They have residue fields
$\mathbb F_{17}$ and $\mathbb F_{289}$,
respectively. The norm of a unit in a
totally ramified extension of degree
$3^d$ reduces to the $3^d$th power
of its residue. This follows either
by taking the determinant of
multiplication on the local integral
algebra and reducing its length
filtration, or by multiplying its
conjugates in a normal closure:
all residue embeddings are the
identity and occur with ramification
multiplicity $3^d$.

For $d=2$ the residue norm image is the
ninth-power subgroup of order thirty-two.
The residue of $\zeta_3$ has order
three, so it is not in this subgroup;
explicitly
$\bar\zeta_3^{\,32}=\bar\zeta_3^2\ne1$.
Any element whose norm has valuation
zero must itself have valuation zero,
since the residue degree is one.
It follows that $\zeta_3$ is not a norm
from $L_{2,17}$.
The same obstruction holds for the
completed multiplicative groups:
it factors through the finite quotient
$\mathbb F_{289}^{\times}/
(\mathbb F_{289}^{\times})^9$.
Thus it also forbids a norm equal to
the completed class of $\zeta_3$.
The complex component of the actual
map $R_3\to R_1$ is $N_{L_2/K}$.
Local restriction of any purported
lift would contradict this calculation,
proving the global failure at $H_3$.

One can see the first lift in the same
finite field. With
$\bar\omega^2+\bar\omega+1=0$,
the element (E62.11) has residue
\[
       \bar w=2+8\bar\omega,\qquad
       \bar w^3=\bar\omega,\qquad
       \bar w^{96}=\bar\omega^2 .
 \tag{E62.13}
\]
Thus it is a primitive ninth root
whose cubic residue is nontrivial.
Every possible first preimage of
$\zeta_3$ encounters the same
second-step obstruction, since the
composite norm has ninth-power
residue image.

\emph{Evaluation of every local level and the limit.}
Principal units in a finite extension
of $\mathbb Q_{17}$ form a pro-seventeen
group. Raising to $3^m$ is an
automorphism on that group, by Hensel's
lemma, so it contributes zero to the
completion for multiplication by
$3^m$. The Teichmuller units in
$F_{d,17}$ have order sixteen and
also contribute zero. The
three-primary Teichmuller units in
$L_{d,17}$ are exactly the same
$\mu_9\subset K_{17}$ at every level:
a totally ramified extension adds
no unramified roots of order prime
to seventeen.
Consequently the actual coordinate
maps are
\[
 \begin{aligned}
 F_{d,17}^{\times,\wedge_3}
       &\xrightarrow{\ \sim\ }\mathbb Z_3,
              &\alpha_d^e&\longleftrightarrow e,\\
 L_{d,17}^{\times,\wedge_3}
       &\xrightarrow{\ \sim\ }\mathbb Z_3\oplus\mathbb Z/9,
              &\beta_d^f\iota^t&\longleftrightarrow(f,t).
 \end{aligned}
 \tag{E62.14}
\]
The uniformizer norms are
$N(\alpha_{d+1})=\alpha_d$ and
$N(\beta_{d+1})=\beta_d$, from
their odd-degree monic polynomials.
The norm of the common root
$\iota$ is $\iota^3$.
This proves the transition
$(e,f,t)\mapsto(e,f,3t)$.
The inverse limit of
$\mathbb Z/9\xleftarrow{3}
\mathbb Z/9\xleftarrow{3}\cdots$
is zero, since any coordinate is
divisible by nine after two further
transitions. The other two coordinates
are unchanged.

The total local norm of the two
uniformizer coordinates is
$17^{e+2f}$.
The ninth root has trivial norm
to $\mathbb Q_{17}$, since the
unramified quadratic norm on it
raises it to $17+1=18$.
The condition for belonging to
$H^1(J_{A,n})$ is therefore exactly
$e+2f\in3\mathbb Z_3$.
This proves (E62.6)--(E62.8).
The continuous inverse-limit
comparison is the actual
cocycle-alignment map of E.46:
the positional module maps are
surjective, and every positional
Tate module has zero invariant
subgroup. Compatible classes have
compatible cocycle representatives,
uniquely up to a compatible
bounding element. All the groups
in (E62.14) are compact, and the
displayed coordinate maps and
their inverses are continuous.
This also proves the stated topology.

Finally the toric map divides both
raw coordinates by three.
Its second class is the Kummer
class of $17^{e+2f}$ divided by
three, giving (E62.9).
The integral map
$(b,f)\mapsto(3b-2f,f)$ is a
bijection from $\mathbb Z_3^2$
to $\mathcal L$, with inverse
$((e+2f)/3,f)$.
This explicitly evaluates every
toric fiber and proves its
surjectivity. The finite-field
certificate checks all residue
norm images used above, the
literal element (E62.13), the
torsion transitions, and the
lattice coordinates modulo
$3,9,27,81$.
\end{proof}
